%%%%%%%%%%%%%%%%%%%%%%%%%%%%%%%%%%%%%%%%%%%%%%%%%%%%%%%%%%%%
\section{SB15 -- RING (Residue Interaction Networks)}

\subsection*{Table of Contents}

\begin{enumerate}
\item Graph representation of proteins
\item Residue interaction networks
\item Nodes and edges
\item Contact maps
\item Non-covalent interaction types
\item Network construction
\item Network topology
\item Centrality measures
\item Functional interpretation of networks
\item Applications of RING
\end{enumerate}


\begin{enumerate}
\item{What is a residue interaction network (RIN)?}
\item{Why can proteins be represented as graphs?}
\item{What do nodes and edges represent in a residue interaction network?}
\item{Which interaction types can define edges in a RIN?}
\item{How is a contact map constructed?}
\item{What information is lost when converting a protein structure into a graph?}
\item{How can centrality measures identify important residues?}
\item{How are residue interaction networks used to study mutations?}
\item{What is the difference between atomic and residue-level interaction networks?}
\item{What are the strengths and weaknesses of the RING approach?}
\end{enumerate}